\documentclass[10pt,a4paper]{amsart}
\usepackage[margin=19mm]{geometry}
\usepackage{amsmath,amssymb,mathtools}
\usepackage[scr=rsfs,cal=euler]{mathalfa}
\usepackage{xfrac}
\usepackage[unicode,hidelinks,bookmarks=false]{hyperref}
\newcommand{\ie}{i.e.\ }
\newcommand{\cf}{cf.\ }
\newcommand{\resp}{respectively\ }
\newcommand{\Subsection}{\subsection}
\providecommand{\nobreakdash}{\nobreak\hbox{-}}
\setlength{\emergencystretch}{2em}
\multlinegap0pt
\usepackage{amssymb}
\usepackage{mathtools}
\newtheorem{proposition}{Proposition}
\renewcommand{\H}{\mathcal H}
\newcommand{\la}{\langle}
\newcommand{\qtq}[1]{\quad\text{#1}\quad}
\newcommand{\ra}{\rangle}
\title{Existence and equicontinuity of solutions\\to the Kirchhoff--Pohozaev wave equation}
\author{Luccas Campos, Rowan Killip, Monica Vi\c{s}an}
\date{Source: arXiv:2608.25047v1 (CC BY 4.0)}
\begin{document}
\maketitle

\noindent\emph{Source and licence.} Existence and equicontinuity of solutions\\to the Kirchhoff--Pohozaev wave equation. Version arXiv:2608.25047v1 (CC BY 4.0), distributed by arXiv under the Creative Commons Attribution 4.0 International licence, http://creativecommons.org/licenses/by/4.0/, as recorded in the arXiv licence metadata for this exact version and shown on the abstract page https://arxiv.org/abs/2608.25047v1. Full licence notice: \texttt{licenses/arxiv-2608.25047-CC-BY-4.0.txt}.\par\smallskip

\noindent\textit{Editorial scope.} Counted unit: the article's proposition P:Poisson (source0.tex lines 383--392). The article's complete original proof of it is delivered with it (source0.tex lines 394--499, one \texttt{proof} environment of 106 lines). No mathematics is written, repaired or re-derived: the statement and the proof are the article's own lines, in the article's own notation, and no retained line is substituted or re-flowed.  Retained uncounted as the source-local context the proof needs: lines 252--254.  Nothing was omitted from any retained span, so the package carries no omission record.  No retained line contains a citation, so the wrapper carries no bibliography and no bibliography entry.  No retained line contains a figure environment, an image-inclusion command or drawing code, so no image is delivered and the package declares no asset dependency.  Editorial formatting only, in addition: an AMS \texttt{amsart} preamble that restates the article's own declarations and macros used by the retained lines, namely the environments \texttt{proposition} on separate counters, together with the standard packages those lines need.  The retained environments print in this document as Proposition 1 in order of appearance, whereas the article prints the same items with the numbering of its own full document.  The complete native source of the article as distributed in the arXiv e-print for this exact version is bundled unchanged as \texttt{sources/208538/source0.tex}.
\begin{equation}\label{KP+}
\tfrac{d}{dt}q = p, \quad \tfrac{d}{dt}p = -c^2(q) A q \qtq{in $\H^{s-1}$ sense with}c(q) := \tfrac{1}{1 + \langle q,A q\rangle} .
\end{equation}

\begin{proposition}\label{P:Poisson} For any $z,w \geq 0$, we have
\begin{equation}\label{Poisson commute}
\{I(w;q,p),I(z;q,p)\}=0 \qtq{on} \H^\infty.
\end{equation}
In particular, as $I(0;q,p) =2H(q,p)$, the quantity $I(z)$ is conserved for all $\H^\infty$ solutions to \eqref{KP+}. More generally,
\begin{align}\label{1153}
\{I(z;q,p),I_k(q,p)\}=0 \qtq{and} \{I_k(q,p),I_\ell(q,p)\}=0
\end{align}
for all $z\geq0$ and integers $k,\ell\geq 0$.
\end{proposition}

\begin{proof}
We decompose $I (z)=I^{(1)}(z)+I^{(2)}(z)+I^{(3)}(z)$, where
\begin{align*}
I^{(1)}(z)&:=\langle p,\tfrac{1}{1+zA}p\rangle, \quad  I^{(2)}(z) := c(q)\langle q,\tfrac{A}{1+zA}q\rangle \qtq{and} \\
I^{(3)}(z)& := -z\Bigl[\langle q,\tfrac{A}{1+zA}q\rangle\langle p,\tfrac{A}{1+zA}p\rangle-\langle p,\tfrac{A}{1+zA}q\rangle^2\Bigr].
\end{align*}
We have 
\begin{align}
\tfrac{\partial}{\partial q} I^{(1)} (z)&= 0 \qtq{and} \tfrac{\partial}{\partial p} I^{(1)} (z)= 2 \tfrac{1}{1+zA}p \label{eq:partial 1} \\
\tfrac{\partial}{\partial q} I^{(2)} (z)&= 2 c(q)  \tfrac{A}{1+zA}q - 2 c(q)^2\langle q, \tfrac{A}{1+zA}q\rangle Aq \qtq{and} \tfrac{\partial}{\partial p} I^{(2)}(z) = 0\label{eq:partial 2}\\
\tfrac{\partial}{\partial q} I^{(3)} (z)&= -2z\langle p, \tfrac{A}{1+zA}p\rangle  \tfrac{A}{1+zA}q +2z \langle p, \tfrac{A}{1+zA}q\rangle  \tfrac{A}{1+zA}p \label{eq:partial 3}\\
\tfrac{\partial}{\partial p} I^{(3)}(z) &= -2z\langle q, \tfrac{A}{1+zA}q\rangle  \tfrac{A}{1+zA}p +2z \langle p, \tfrac{A}{1+zA}q\rangle  \tfrac{A}{1+zA}q. \label{eq:partial 4}
\end{align}

By \eqref{eq:partial 1} and \eqref{eq:partial 2}, it is easy to see that
\begin{align}\label{eq:poisson_commutativity 1}
\{I^{(1)}(w),I ^{(1)}(z) \}=0 \qtq{and} \{I^{(2)}(w),I ^{(2)}(z) \}=0.
\end{align}

Using \eqref{eq:partial 1} and \eqref{eq:partial 3}, we find
\begin{align*}
&\{I^{(1)}(w),I^{(3)}(z)\} + \{I^{(3)}(w),I^{(1)}(z)\}\\ 
&=-\bigl\langle \tfrac{\partial}{\partial p} I^{(1)}(w),\tfrac{\partial}{\partial q} I^{(3)}(z) \bigr\rangle
	+ \bigl\langle \tfrac{\partial}{\partial q} I^{(3)}(w),\tfrac{\partial}{\partial p} I^{(1)}(z) \bigr\rangle\\
&= 4\la p,\tfrac{A}{(1+zA)(1+wA)}q\ra\la p, \tfrac{zA}{1+z A}p\ra- 4\la p,\tfrac{A}{(1+zA)(1+wA)}p\ra\la p, \tfrac{zA}{1+zA}q\ra\\
&\quad-4\la p,\tfrac{A}{(1+zA)(1+wA)}q\ra\la p, \tfrac{wA}{1+wA}p\ra+4\la p,\tfrac{A}{(1+zA)(1+wA)}p\ra\la p, \tfrac{wA}{1+wA}q\ra.
\end{align*}
Employing the identity
\begin{equation}\label{id}
 \tfrac{zA}{1+z A}- \tfrac{wA}{1+w A} = (z-w) \tfrac{A}{(1+z A)(1+wA)},
\end{equation}
this simplifies to 
\begin{align}\label{eq:poisson_commutativity 2}
\{I^{(1)}(w),I^{(3)}(z)\} &+ \{I^{(3)}(w),I^{(1)}(z)\} \\ 
&=4(z-w)\la p,\tfrac{A}{(1+zA)(1+wA)}q\ra\la p, \tfrac{A}{(1+z A)(1+wA)} p\ra\notag\\
&\quad-4(z-w)\la p,\tfrac{A}{(1+zA)(1+wA)}p\ra\la p, \tfrac{A}{(1+z A)(1+wA)} q\ra\notag\\
&=0.\notag
\end{align}

Next we show that
\begin{align}\label{eq:poisson_commutativity 4}
\{I^{(3)}(w),I ^{(3)}(z) \}=0.
\end{align}
A straightforward computation employing \eqref{eq:partial 3} and \eqref{eq:partial 4} yields
\begin{align*}
\{I^{(3)}&(w),I ^{(3)}(z) \}\\
&=4zw \la p, \tfrac{A^2}{(1+zA)(1+wA)}q\ra\Bigl[  \la q, \tfrac{A}{1+zA}q\ra\la p, \tfrac{A}{1+wA}p\ra -  \la q, \tfrac{A}{1+wA}q\ra\la p, \tfrac{A}{1+zA}p\ra\Bigr]\\
&\quad-4zw\la p, \tfrac{A^2}{(1+zA)(1+wA)}p\ra\Bigl[  \la q, \tfrac{A}{1+zA}q\ra\la p, \tfrac{A}{1+wA}q\ra -  \la q, \tfrac{A}{1+wA}q\ra\la p, \tfrac{A}{1+zA}q\ra\Bigr]\\
&\quad-4zw\la q, \tfrac{A^2}{(1+zA)(1+wA)}q\ra\Bigl[  \la p, \tfrac{A}{1+zA}q\ra\la p, \tfrac{A}{1+wA}p\ra -  \la p, \tfrac{A}{1+wA}q\ra\la p, \tfrac{A}{1+zA}p\ra\Bigr].
\end{align*}
To deduce the claim \eqref{eq:poisson_commutativity 4}, we apply the following identity in each of the square brackets above:
\begin{align}\label{id 2}
 \la f, \tfrac{A}{1+zA}g\ra\la \phi, \tfrac{A}{1+wA}\psi\ra &-  \la f, \tfrac{A}{1+wA}g\ra\la \phi, \tfrac{A}{1+zA}\psi\ra\notag\\
 &=(z-w) \la f, \tfrac{A}{(1+zA)(1+wA)}g\ra \la \phi, \tfrac{A^2}{(1+zA)(1+wA)}\psi\ra\\
 &\quad - (z-w) \la f, \tfrac{A^2}{(1+zA)(1+wA)}g\ra \la \phi, \tfrac{A}{(1+zA)(1+wA)}\psi\ra .\notag
\end{align}
More precisely, we take $f=g=q$ and $\phi=\psi=p$ in the first term, $f=g=\psi=q$ and $\phi=p$ in the second, and $f=\phi=\psi=p$ and $g=q$ in the third.

In view of \eqref{eq:poisson_commutativity 1}, \eqref{eq:poisson_commutativity 2}, and \eqref{eq:poisson_commutativity 4}, claim \eqref{Poisson commute} will follow once we show
\begin{align}\label{eq:poisson_commutativity 3}
\{I^{(1)}(w)+I^{(3)}(w),I ^{(2)}(z) \} +	\{I^{(2)}(w),I ^{(1)}(z)+I ^{(3)}(z) \} =0.
\end{align}
To this end, we use \eqref{eq:partial 2} and \eqref{eq:partial 4} to compute
\begin{align*}
&\{I^{(2)}(w),I^{(3)}(z) \} \\
&= \bigl\langle \tfrac{\partial}{\partial q} I^{(2)}(w),\tfrac{\partial}{\partial p} I^{(3)}(z) \bigr\rangle\\
&=-4c(q)\la q, \tfrac{zA}{1+zA}q\ra \la p, \tfrac{A^2}{(1+z A)(1+wA)} q\ra +4c(q)\la p, \tfrac{zA}{1+zA}q\ra \la q, \tfrac{A^2}{(1+z A)(1+wA)} q\ra\\
&\quad +4 c(q)^2 \la q, \tfrac{A}{1+wA}q\ra\la q, \tfrac{A}{1+zA}q\ra\la p, \tfrac{zA^2}{1+zA}q\ra \\
&\quad-4 c(q)^2 \la q, \tfrac{A}{1+wA}q\ra\la p, \tfrac{A}{1+zA}q\ra\la q, \tfrac{zA^2}{1+zA}q\ra\\
&=-4c(q)\la q, \tfrac{zA}{1+zA}q\ra \la p, \tfrac{A^2}{(1+z A)(1+wA)} q\ra +4c(q)\la p, \tfrac{zA}{1+zA}q\ra \la q, \tfrac{A^2}{(1+z A)(1+wA)} q\ra\\
&\quad +4 c(q)^2 \la q, \tfrac{A}{1+wA}q\ra\la q, \tfrac{A}{1+zA}q\ra\la p, \tfrac{zA^2}{1+zA}q\ra \\
&\quad-4 c(q)^2 \la q, \tfrac{A}{1+wA}q\ra\la p, \tfrac{A}{1+zA}q\ra\Bigl[ \la q, Aq\ra-  \la q, \tfrac{A}{1+zA}q\ra\Bigr] .
\end{align*}
From this (both as written and after interchanging the $z$ and $w$ variables) and the identity \eqref{id}, we obtain
\begin{align}\label{441}
\{I^{(2)}(w),I ^{(3)}(z) \} & +\{I^{(3)}(w),I ^{(2)}(z) \} \notag\\
&=-4(z-w)c(q)\la q, \tfrac{A}{(1+z A)(1+wA)}q\ra \la p, \tfrac{A^2}{(1+z A)(1+wA)} q\ra\notag\\
&\quad+4(z-w)c(q)\la p, \tfrac{A}{(1+z A)(1+wA)}q\ra \la q, \tfrac{A^2}{(1+z A)(1+wA)} q\ra\notag\\
&\quad +4 (z-w) c(q)^2 \la q, \tfrac{A}{1+wA}q\ra\la q, \tfrac{A}{1+zA}q\ra\la p, \tfrac{A^2}{(1+zA)(1+wA)}q\ra\notag\\
&\quad- 4(z-w) c(q)^2 \la q, \tfrac{A}{1+wA}q\ra\la q, \tfrac{A}{1+zA}q\ra\la p, \tfrac{A^2}{(1+zA)(1+wA)}q\ra \notag\\
&\quad + 4c(q)^2 \la q,Aq\ra \Bigl[ \la q, \tfrac{A}{1+zA}q\ra \la p, \tfrac{A}{1+wA}q\ra - 
	\la q, \tfrac{A}{1+wA}q\ra \la p, \tfrac{A}{1+zA}q\ra  \Bigr] \notag\\
&=-4(z-w)c(q)\la q, \tfrac{A}{(1+z A)(1+wA)}q\ra \la p, \tfrac{A^2}{(1+z A)(1+wA)} q\ra\\
&\quad+4(z-w)c(q)\la p, \tfrac{A}{(1+z A)(1+wA)}q\ra \la q, \tfrac{A^2}{(1+z A)(1+wA)} q\ra\notag\\
&\quad + 4c(q)^2 \la q,Aq\ra \Bigl[ \la q, \tfrac{A}{1+zA}q\ra \la p, \tfrac{A}{1+wA}q\ra - 
	\la q, \tfrac{A}{1+wA}q\ra \la p, \tfrac{A}{1+zA}q\ra  \Bigr]. \notag
\end{align}


Using \eqref{eq:partial 1} and \eqref{eq:partial 2}, we compute
\begin{align}\label{440}
\{I^{(1)}(w),&\, I ^{(2)}(z) \}+ \{I^{(2)}(w),I ^{(1)}(z) \} \notag\\
&= -\bigl\langle \tfrac{\partial}{\partial p} I^{(1)}(w),\tfrac{\partial}{\partial q} I^{(2)}(z) \bigr\rangle
	+ \bigl\langle \tfrac{\partial}{\partial q} I^{(2)}(w),\tfrac{\partial}{\partial p} I^{(1)}(z) \bigr\rangle\\
&=4c(q)^2 \Bigl[ \la q,\tfrac{A}{1+zA}q\ra \la p, \tfrac{A}{1+wA}q\ra  -  \la q, \tfrac{A}{1+wA}q\ra \la p,\tfrac{A}{1+zA}q\ra \Bigr], \notag
\end{align}
which we then combine with \eqref{441} to obtain
\begin{align*}
\{I^{(1)}&(w)+I^{(3)}(w),I ^{(2)}(z) \} +	\{I^{(2)}(w),I ^{(1)}(z)+I ^{(3)}(z) \} \\
&=-4(z-w)c(q)\la q, \tfrac{A}{(1+z A)(1+wA)}q\ra \la p, \tfrac{A^2}{(1+z A)(1+wA)} q\ra\\
&\quad+4(z-w)c(q)\la p, \tfrac{A}{(1+z A)(1+wA)}q\ra \la q, \tfrac{A^2}{(1+z A)(1+wA)} q\ra\notag\\
&\quad + 4c(q) \Bigl[ \la q, \tfrac{A}{1+zA}q\ra \la p, \tfrac{A}{1+wA}q\ra - 
	\la q, \tfrac{A}{1+wA}q\ra \la p, \tfrac{A}{1+zA}q\ra  \Bigr]. \notag
\end{align*}
Lastly, we apply the identity \eqref{id 2} with $f=g=\psi= q$ and $\phi=p$ to deduce the result \eqref{eq:poisson_commutativity 3}.  The claims in \eqref{1153} follow by examining Taylor coefficients about $w=0$ and $z=0$.
\end{proof}

\end{document}
